\chapter{Skrip Utama Orkestrasi dan Algoritma Hibrida Kuantum}

Pada bagian ini, dilampirkan kumpulan skrip pemrograman utama yang mengatur alur logika eksekusi algoritma GT-VQE secara terintegrasi.

\begin{lstlisting}[language=Python, caption={Skrip utama orkestrasi keseluruhan parameter dan eksekusi pada sistem N=4 (\texttt{main\_N4.py})}, label={lst:main-N4}]
import warnings
import pandas as pd
from data_downloader import download_market_data
from backtest_runner import run_backtest
from report_generator import generate_report
from plot_generator import plot_all
from config import get_config, clean_workspace

warnings.filterwarnings('ignore')

tickers = ['BBCA.JK', 'TLKM.JK', 'SMGR.JK', 'ADRO.JK']
N = len(tickers)

# 0. Bersihkan Workspace
clean_workspace(N)

config = get_config(tickers)

suffix = f"_N{N}"
filenames = {
    'daily_prices_csv': f'harga_harian_saham{suffix}.csv',
    'report_txt': f'laporan_backtest{suffix}.txt',
    'circuit_png': f'rangkaian_kuantum{suffix}.png',
    'convergence_png': f'grafik_konvergensi_detail{suffix}.png',
    'backtest_vqe_png': f'hasil_backtest_vqe{suffix}.png',
    'backtest_nash_png': f'hasil_backtest_nash{suffix}.png',
    'depth_png': f'grafik_depth_per_window{suffix}.png'
}

data_clean, benchmark_data, benchmark_rets = download_market_data(
    config['tickers'], config['benchmark_ticker'], 
    config['start_date'], config['end_date'], config['start_bt_date']
)
data_clean.to_csv(filenames['daily_prices_csv'])
print(f"Data harga harian telah diekspor ke '{filenames['daily_prices_csv']}'.")

results = run_backtest(data_clean, config['tickers'], config)
value_benchmark_idx = generate_report(results, data_clean, benchmark_data, benchmark_rets, config['tickers'], config, filenames['report_txt'])
plot_all(results, data_clean, value_benchmark_idx, config, filenames)
print(f"\nSeluruh proses N={N} selesai.")
\end{lstlisting}

\begin{lstlisting}[language=Python, caption={Skrip inti loop VQE adaptif (\texttt{run\_vqe\_adaptive.py})}, label={lst:run-vqe-adaptive}]
import numpy as np
import pandas as pd
import os
import pennylane as qml
from compute_metrics import calculate_entanglement_entropy

def run_vqe_adaptive(H_obj, H_pen, n_qubits, curr_date, ne_bitstring=None, K=2, target_penalty=5.0,
                     max_depth=4, maxiter=100, max_total_iter=500,
                     batch_size=10, conv_window=4, conv_tol=1e-3,
                     best_a_base=0.1, seed=42, file_config=None):
    """
    Menjalankan VQE dengan Penalty Annealing.
    H_total(t) = H_obj + penalty(t) * H_pen
    """
    if file_config is None:
        file_config = {
            'riwayat_iterasi': 'riwayat_iterasi_vqe.csv',
            'theta_final': 'theta_final_all_depths.csv',
            'depth_vs_energi': 'hasil_depth_vs_energi.csv'
        }
        
    dev = qml.device("default.qubit", wires=n_qubits)
    rng = np.random.default_rng(seed)

    # --- RESET RIWAYAT ITERASI ---
    if os.path.exists(file_config['riwayat_iterasi']):
        os.remove(file_config['riwayat_iterasi'])

    def make_circuit(depth):
        def Ansatz(params):
            w = params.reshape((depth + 1, n_qubits, 2))
            for layer in range(depth + 1):
                for q in range(n_qubits):
                    qml.RY(w[layer, q, 0], wires=q)
                    qml.RZ(w[layer, q, 1], wires=q)
                if layer < depth:
                    for q in range(n_qubits - 1):
                        qml.CNOT(wires=[q, q + 1])
                    qml.CNOT(wires=[n_qubits - 1, 0])

        @qml.qnode(dev)
        def cost_circuit(params, p_scale=1.0):
            Ansatz(params)
            # H_total dinamis: H_obj + scale * H_pen
            return qml.expval(H_obj + p_scale * H_pen)

        @qml.qnode(dev)
        def prob_circuit(params):
            Ansatz(params)
            return qml.probs(wires=range(n_qubits))
            
        @qml.qnode(dev)
        def state_circuit(params):
            Ansatz(params)
            return qml.state()

        return cost_circuit, prob_circuit, state_circuit

    def run_spsa(cost_circuit, state_circuit, n_params, init_params=None, a_base=0.1, c_base=0.1, depth_label=1, target_p=5.0):
        a, c      = a_base, c_base
        A_coeff   = maxiter * 0.1
        alpha_exp = 0.602
        gamma_exp = 0.101

        if init_params is not None:
            params = init_params.copy()
        else:
            params = rng.uniform(0, 2 * np.pi, n_params)

        energy_history = []
        entropy_history = []
        grad_var_history = []
        total_iters    = 0
        
        while total_iters < max_total_iter:
            # --- PENALTY ANNEALING TERKALIBRASI ---
            # Pinalti mencapai 100% pada 70% total budget iterasi.
            # Sisa 30% digunakan untuk "fine-tuning" dengan pinalti konstan.
            anneal_progress = min(1.0, total_iters / (0.7 * max_total_iter))
            p_scale = target_p * (0.1 + 0.9 * anneal_progress)
            
            current_grad_vars = []
            for _ in range(batch_size):
                k     = total_iters
                a_k   = a / (A_coeff + k + 1) ** alpha_exp
                c_k   = c / (k + 1)           ** gamma_exp
                delta = 2 * rng.integers(0, 2, size=n_params) - 1
                
                cost_plus  = float(cost_circuit(params + c_k * delta, p_scale=p_scale))
                cost_minus = float(cost_circuit(params - c_k * delta, p_scale=p_scale))
                grad       = (cost_plus - cost_minus) / (2 * c_k * delta)
                
                # [Phase 3] Gradient Clipping
                grad       = np.clip(grad, -1.0, 1.0)
                
                params     = params - a_k * grad
                current_grad_vars.append(np.var(grad))
                total_iters += 1

            current_energy = float(cost_circuit(params, p_scale=p_scale))
            energy_history.append(current_energy)
            grad_var_history.append(np.mean(current_grad_vars))
            
            # Hitung Entropi
            try:
                current_psi = state_circuit(params)
                current_entropy = calculate_entanglement_entropy(current_psi)
                entropy_history.append(current_entropy)
            except:
                entropy_history.append(0.0)
            
            iter_df = pd.DataFrame({
                'Depth': [depth_label], 
                'Iteration': [total_iters], 
                'Energy': [current_energy],
                'Entropy': [entropy_history[-1]],
                'Grad_Var': [grad_var_history[-1]],
                'Penalty_Scale': [p_scale]
            })
            iter_df.to_csv(file_config['riwayat_iterasi'], mode='a', header=not os.path.exists(file_config['riwayat_iterasi']), index=False)

            if total_iters >= maxiter and len(energy_history) >= conv_window:
                # [Phase 4] Entropy Protection & Annealing Protection
                # Jangan berhenti jika:
                # 1. Entropi masih tinggi (mencegah superposisi macet)
                # 2. Annealing pinalti belum mencapai 100%
                if (n_qubits <= 4 and entropy_history[-1] > 0.1) or (anneal_progress < 1.0):
                    continue 

                E_old, E_now = energy_history[-conv_window], energy_history[-1]
                if abs(E_old - E_now) / (abs(E_old) + 1e-12) < conv_tol:
                    break
        return params, energy_history[-1], energy_history, entropy_history, grad_var_history, total_iters

    best_energy    = np.inf
    best_params, best_depth, best_history, best_ent_hist, best_grad_var_hist = None, 1, [], [], []
    depth_energies = []
    prev_params    = None
    all_theta_data = []

    for depth in range(1, max_depth + 1):
        n_params = n_qubits * 2 * (depth + 1)
        cost_fn, prob_fn, state_fn = make_circuit(depth)

        if prev_params is not None and len(prev_params) < n_params:
            # Warm-start yang diperbaiki: gunakan distribusi normal sangat kecil 
            # agar layer baru mendekati identitas, menjaga stabilitas energi.
            old_w = prev_params.reshape((depth, n_qubits, 2))     # (depth, n, 2)
            new_layer = rng.normal(0, 1e-4, (1, n_qubits, 2))    # layer baru (dekat identitas)
            # Sisipkan layer baru sebelum output layer terakhir
            new_w = np.concatenate([old_w[:-1], new_layer, old_w[-1:]], axis=0)  # (depth+1, n, 2)
            
            # [Phase 2] Asymmetric Jitter: Pemutus simetri agar tidak terjebak superposisi
            asymmetric_jitter = rng.normal(0, 1e-4, n_params) * np.linspace(1, 1.5, n_params)
            init_p = new_w.flatten() + asymmetric_jitter
        else:
            if ne_bitstring is not None:
                init_w = np.zeros((depth + 1, n_qubits, 2))
                for q in range(n_qubits):
                    init_w[0, q, 0] = np.pi if int(ne_bitstring[q]) == 1 else 0.0
                
                # [Phase 2] Asymmetric Jitter pada inisialisasi Nash
                asymmetric_jitter = rng.normal(0, 1e-3, n_params) * np.linspace(1, 1.5, n_params)
                init_p = init_w.flatten() + asymmetric_jitter
            else:
                init_p = rng.uniform(0, 2 * np.pi, n_params)

        # Skala learning rate dinamis: semakin dalam sirkuit, langkah semakin kecil (presisi)
        current_a_base = best_a_base / np.sqrt(depth)
        
        params, energy, e_hist, ent_hist, g_var_hist, n_iters = run_spsa(
            cost_fn, state_fn, n_params, init_params=init_p, 
            a_base=current_a_base, c_base=0.1/np.sqrt(depth), 
            depth_label=depth, target_p=target_penalty
        )
        print(f"    [Depth {depth}] Konvergen dalam {n_iters} iterasi | E = {energy:.6f}")
        depth_energies.append((depth, energy, n_iters))
        prev_params = params

        # --- KUMPULKAN THETA (Poin 1) ---
        for idx, val in enumerate(params):
            all_theta_data.append({'Depth': depth, 'Theta_Index': idx, 'Theta_Value': val})

        if energy < best_energy:
            best_energy, best_params, best_depth, best_history, best_ent_hist, best_grad_var_hist = energy, params, depth, e_hist, ent_hist, g_var_hist
            
    # --- EKSPOR THETA GABUNGAN (Poin 1) ---
    pd.DataFrame(all_theta_data).to_csv(file_config['theta_final'], index=False)

    # --- EKSPOR DEPTH VS ENERGI (Poin 3) ---
    depth_df = pd.DataFrame(depth_energies, columns=['Depth', 'Energy', 'Iterations'])
    depth_df.insert(0, 'Date', curr_date.date())
    depth_df.to_csv(file_config['depth_vs_energi'], mode='a', header=not os.path.exists(file_config['depth_vs_energi']), index=False)

    _, prob_fn, _ = make_circuit(best_depth)
    probs = prob_fn(best_params)
    sorted_indices = np.argsort(probs)[::-1]
    best_bitstring = None
    for idx in sorted_indices:
        bs = format(idx, f'0{n_qubits}b')
        if bs.count('1') == K:
            best_bitstring = bs
            break
    if best_bitstring is None:
        top_k = np.argsort(probs)[-K:]
        bs_list = list('0' * n_qubits)
        for idx in top_k: bs_list[idx] = '1'
        best_bitstring = ''.join(bs_list)

    return [i for i, bit in enumerate(best_bitstring) if bit == '1'], best_depth, best_energy, best_history, best_ent_hist, best_grad_var_hist, depth_energies, probs, best_params
\end{lstlisting}

\begin{lstlisting}[language=Python, caption={Skrip eksekusi langkah strategi yang memadukan teori permainan klasik dengan VQE kuantum (\texttt{run\_strategy\_step.py})}, label={lst:run-strategy-step}]
import numpy as np
import pandas as pd
import os
from compute_endogenous_lambda import compute_endogenous_lambda
from find_nash_sbr import find_nash_sbr
from brute_force_validator import run_brute_force_window
from build_hamiltonian_total import build_hamiltonian_total
from find_optimal_lr_spsa import find_optimal_lr_spsa
from run_vqe_adaptive import run_vqe_adaptive
from calc_simple_return import calculate_simple_return
from calc_log_return import calculate_log_return
from calc_expected_returns import calculate_expected_simple_return, calculate_expected_log_return_with_drag
from calc_expected_returns import calculate_expected_simple_return, calculate_expected_log_return_with_drag

def run_strategy_step(lookback_data, tickers, curr_date, K=2, penalty_A=5.0,
                      max_depth=6, maxiter=100,
                      max_total_iter=500, batch_size=10,
                      conv_window=4, conv_tol=1e-3,
                      use_warm_start=True,
                      file_config=None):
    """
    Eksekusi satu periode pembelajaran dengan pencatatan tanggal.
    """
    if file_config is None:
        file_config = {
            'metrics': 'metrik_return_dan_lambda.csv',
            'bias_h': 'bias_h_total.csv',
            'interaksi_J': 'interaksi_J_total.csv',
            'pencarian_lr': 'hasil_pencarian_lr.csv',
            'parameter_pendamping': 'parameter_pendamping.csv'
        }
    
    n_assets = len(tickers)
    # ... rest of calculations ...
    # Simple Return dihitung manual via fungsi baru
    simple_rets = calculate_simple_return(lookback_data)
    # Log Return dihitung manual via fungsi baru
    log_rets  = calculate_log_return(lookback_data)
    
    binary_st = (log_rets <= 0).astype(int)   # 1 = turun, 0 = naik

    # --- PERHITUNGAN EXPECTED RETURN (Rumus Volatility Drag) ---
    mu_R_daily = calculate_expected_simple_return(simple_rets)
    # Variance log return sebagai dasar volatility drag
    var_r_daily = log_rets.var() 
    mu_r_daily = calculate_expected_log_return_with_drag(mu_R_daily, var_r_daily)

    # Menggunakan return periode (126 hari)
    mu_simple_period = mu_R_daily.values * 126
    mu_log_period = mu_r_daily.values * 126
    
    sigma_log = log_rets.std().values
    sigma_period_log = sigma_log * np.sqrt(126)
    
    # Matriks kovariansi tetap berbasis log return
    sigma_period_matrix = log_rets.cov().values * 126

    # gamma disini mewakili degree of risk-aversion, menggunakan mu_log_period (dengan drag) dan sigma_period_log
    gamma = compute_endogenous_lambda(mu_log_period, sigma_period_log)
    lam = penalty_A # penalty lambda untuk pembatas kardinalitas
    metrics_df = pd.DataFrame({
        'Date': [curr_date.date()] * n_assets,
        'Ticker': tickers,
        'Mu_Simple_Period': mu_simple_period,
        'Mu_Log_Period': mu_log_period,
        'Sigma_Log': sigma_log,
        'Sigma_Period_Log': sigma_period_log,
        'Lambda_RiskAversion': [gamma] * n_assets
    })
    
    metrics_df.to_csv(file_config['metrics'], mode='a', header=not os.path.exists(file_config['metrics']), index=False)

    # --- KONSTRUKSI MATRIKS Q (QUBO) TERPISAH ---
    # 1. Komponen Objektif (Return & Risk)
    Q_diag_obj = np.zeros(n_assets)
    Q_off_obj  = np.zeros((n_assets, n_assets))
    K_sq = K**2
    
    for i in range(n_assets):
        Q_diag_obj[i] = (gamma * sigma_period_matrix[i, i]) / (2.0 * K_sq) - (mu_simple_period[i] / K)
        for j in range(i + 1, n_assets):
            Q_val = (gamma * sigma_period_matrix[i, j]) / (2.0 * K_sq)
            Q_off_obj[i, j] = Q_val
            Q_off_obj[j, i] = Q_val

    # 2. Komponen Pinalti (Constraint)
    Q_diag_pen = np.zeros(n_assets)
    Q_off_pen  = np.zeros((n_assets, n_assets))
    for i in range(n_assets):
        Q_diag_pen[i] = (1.0 - 2.0 * K)
        for j in range(i + 1, n_assets):
            Q_off_pen[i, j] = 1.0
            Q_off_pen[j, i] = 1.0

    # --- KONSTRUKSI PARAMETER ISING TERPISAH ---
    # Fungsi pembantu untuk konversi Q ke h, J, C
    def qubo_to_ising(Q_diag, Q_off, constant=0.0):
        h = np.zeros(n_assets)
        J = np.zeros((n_assets, n_assets))
        for i in range(n_assets):
            sum_Q_ij_half = 0.0
            for j in range(n_assets):
                if i != j:
                    sum_Q_ij_half += Q_off[i, j] / 2.0
                    if i < j:
                        J[i, j] = Q_off[i, j] / 2.0
                        J[j, i] = Q_off[i, j] / 2.0
            h[i] = -((Q_diag[i] / 2.0) + sum_Q_ij_half)
        
        sum_Q_ii_half = np.sum(Q_diag) / 2.0
        sum_Q_ij_half_total = 0.0
        for i in range(n_assets):
            for j in range(i + 1, n_assets):
                sum_Q_ij_half_total += Q_off[i, j] / 2.0
        C = sum_Q_ii_half + sum_Q_ij_half_total + constant
        return h, J, C

    h_obj, J_obj, C_obj = qubo_to_ising(Q_diag_obj, Q_off_obj, 0.0)
    h_pen, J_pen, C_pen = qubo_to_ising(Q_diag_pen, Q_off_pen, float(K_sq))

    # --- EKSPOR BIAS & INTERAKSI TERPISAH ---
    h_df = pd.DataFrame({
        'Date': [curr_date.date()] * n_assets, 
        'Ticker': tickers, 
        'Bias_h_Obj': h_obj, 
        'Bias_h_Pen': h_pen
    })
    h_df.to_csv(file_config['bias_h'], mode='a', header=not os.path.exists(file_config['bias_h']), index=False)
    
    J_list = []
    for i in range(n_assets):
        for j in range(i + 1, n_assets):
            J_list.append({
                'Date': curr_date.date(), 
                'Ticker_i': tickers[i], 
                'Ticker_j': tickers[j], 
                'Interaction_J_Obj': J_obj[i, j],
                'Interaction_J_Pen': J_pen[i, j]
            })
    J_df = pd.DataFrame(J_list)
    J_df.to_csv(file_config['interaksi_J'], mode='a', header=not os.path.exists(file_config['interaksi_J']), index=False)

    # --- EKSPOR KONSTANTA C ISING TERPISAH ---
    c_df = pd.DataFrame({
        'Date': [curr_date.date()],
        'C_Obj': [C_obj],
        'C_Pen': [C_pen],
        'Penalty_A_Target': [penalty_A]
    })
    c_df.to_csv(file_config['parameter_pendamping'], mode='a', header=not os.path.exists(file_config['parameter_pendamping']), index=False)
    
    # Membangun Hamiltonian terpisah (Diperlukan untuk VQE & LR Finder)
    H_obj = build_hamiltonian_total(h_obj, J_obj, n_assets, offset=C_obj)
    H_pen = build_hamiltonian_total(h_pen, J_pen, n_assets, offset=C_pen)
    
    # [Tugas 4] Pencarian Nash Equilibrium (Hanya jika Warm-Start aktif)
    if use_warm_start:
        ne_bitstring, ne_utility = find_nash_sbr(mu_simple_period, sigma_period_matrix, gamma, curr_date, N=n_assets, K=K, history_file=file_config.get('nash_history', 'riwayat_nash_sbr.csv'))
        print(f"    [Nash Eq] Warm-start aktif. Bitstring: {ne_bitstring} | Utility: {ne_utility:.6f}")
        vqe_init_bs = ne_bitstring
    else:
        ne_bitstring, ne_utility = "N/A", 0.0
        vqe_init_bs = "0" * n_assets # Mulai dari |0...0>
        print("    [VQE] Warm-start dinonaktifkan. Memulai dari $|0...0\\rangle$")

    # [Validasi] Jalankan Brute Force Validation per window
    # Menggunakan Hamiltonian gabungan (Total) untuk pembanding energi VQE
    h_total = h_obj + penalty_A * h_pen
    J_total = J_obj + penalty_A * J_pen
    C_total = C_obj + penalty_A * C_pen
    
    print(f"    [Brute Force] Validasi global minimum untuk window {curr_date.date()}...")
    bf_bs, bf_e = run_brute_force_window(curr_date, n_assets, h_total, J_total, C_total, K, file_config)
    print(f"    [Brute Force] Global Min: {bf_bs} | Energy: {bf_e:.6f}")

    # [LR Finder] Menggunakan H gabungan dengan pinalti target bulan ini
    H_target = H_obj + penalty_A * H_pen
    print("    [LR Finder] Mencari Learning Rate optimal...")
    best_a, test_a_values, final_energies = find_optimal_lr_spsa(H_target, n_assets, curr_date, ne_bitstring=vqe_init_bs, K=K, test_iters=30, file_config=file_config)
    print(f"    [LR Finder] Base Learning Rate terpilih: {best_a:.4f}")

    # Optimasi VQE dengan Penalty Annealing
    selected_indices, depth_used, energy_final, best_history, best_ent_hist, best_gvar_hist, depth_energies, winning_probs, best_params = run_vqe_adaptive(
        H_obj, H_pen, n_assets, curr_date, ne_bitstring=vqe_init_bs, K=K, target_penalty=penalty_A,
        max_depth=max_depth, maxiter=maxiter,
        max_total_iter=max_total_iter, batch_size=batch_size,
        conv_window=conv_window, conv_tol=conv_tol, 
        best_a_base=best_a, file_config=file_config
    )

    lr_data = (test_a_values, final_energies, best_a)
    return selected_indices, depth_used, energy_final, best_history, best_ent_hist, best_gvar_hist, depth_energies, lr_data, ne_bitstring, ne_utility, winning_probs, best_params
\end{lstlisting}

